\part{Функциональный анализ.}
\chapter{Линейные нормированные пространства.}
\begin{definition}
    Пусть $X$ -- вещественное линейное пространство.
    Отображение $\|\cdot\|\colon X \to \R$ называется нормой, если
    \begin{enumerate}
        \item $\forall x \in X \hookrightarrow \|x\| \geq 0 \wedge \|x\| = 0 \Leftrightarrow x = 0$.
        \item $\forall x \in X \ \forall \alpha \in \R \hookrightarrow \|\alpha x\| = |\alpha|\|x\|$.
        \item $\forall x, y \in X \hookrightarrow \|x + y\| \leq \|x\| + \|y\|$.
    \end{enumerate}

    Пара $(X, \|\cdot\|)$ называется нормированным пространством.
\end{definition}
\begin{note}
    Всякая норма порождает метрику как
    \[
        \rho(x, y) = \|x - y\|.
    \]
\end{note}
\begin{example}
    Пусть $(X, \|\cdot\|)$ -- нормированное пространство.
    Тогда если $B_R(x) \supset B_r(y)$, то $R \geq r$.

    Если вдруг $x = y$, то для всякого вектора единичной нормы $z\colon$ вектор $x + rz \in B_r(y)$, и, в частности лежит в $B_R(x)$, то есть
    \[
        \|rz\| \leq R \Rightarrow r \leq R.
    \]

    Если же вдруг так получилось, что $x \neq y$, то рассмотрим вектор
    \[
        y + r \frac{x - y}{\|x - y\|} \in B_r(y) \subset B_R(x).
    \]
    Тогда
    \[
        \left\|x - y - r \frac{x - y}{\|x - y\|}\right\| \leq R \Rightarrow \left|1 - \frac{r}{\|x - y\|}\right|\|x - y\| \leq R \Rightarrow |\|x - y\| - r| \leq R.
    \]
    Аналогичное неравенство верно и для $y - r \frac{x - y}{\|x - y\|}\colon$
    \[
        |\|x - y\| + r| \leq R.
    \]
    Тогда
    \[
        2r = |r - \|x - y\| + r + \|x - y\|| \leq |r - \|x - y\|| + |r + \|x - y\|| \leq 2R,
    \]
    что и требовалось показать.
\end{example}
\begin{example}
    В общих метрических пространствах это утверждение неверно: пусть $X = [0, 1] \cup \{2\}$ с евклидовой метрикой.
    Тогда
    \[
        B_{1.5}(0) \subset B_{1}(1).
    \]
    И шар большего радиуса содержится в шаре меньшего радиуса.
\end{example}
\begin{definition}
    Пусть $(X, \|\cdot\|)$ -- нормированное пространство.
    Будем говорить, что оно является банаховым, если оно полно по метрике, порождённой нормой.
\end{definition}
\begin{theorem}
    Пусть $(X, \|\cdot\|)$ -- банахово пространство.
    Тогда для всякой последовательности вложенных шаров $\{B_{r_n}(x_n)\}$ верно, что
    \[
        \bigcap_{n = 1}^{\infty} B_{r_n}(x_n) \neq \emptyset.
    \]
\end{theorem}
\begin{proof}
    В силу вложенности шаров $r_m \leq r_n$, при $m \leq n$.

    Поскольку все радиусы шаров -- неотрицательны, то существует $r = \lim r_n = \inf r_n \geq 0$.

    Пусть $m \geq n$.
    Тогда $B_{r_m}(x_n) \subset B_{r_n}(x_n)$, а значит $\|x_n - x_m\| \leq r_n - r_m$.
    Следовательно, последовательность центров шаров -- фундаментальна, и сходится к некоторой точке $x$, которая и будет лежать в пересечении.
\end{proof}